% !PRINTED 160
\weberpage{160}
%<<CONT
the highest power of $y$ drops out of the final equation. One
could then save the agreement of the number with (5) only by
the arbitrary addition of infinite roots. Multiple roots, too, give
cause for misgivings. These drawbacks are met in part by taking
as a basis the \gesperrt{homogeneous} form of the equations, a form in
which the elimination problem presents itself especially in
geometry.

For if the function $f(x)$ has arisen from a homogeneous
function of order $n$ of the three variables $x, y, z$ (a ternary
form) by arranging according to powers of $x$, then the
coefficients $a_0, a_1, a_2 \ldots a_n$ are homogeneous forms of the
two variables $y, z$ of the degree indicated by the index, so
that $a_0$ is a constant, $a_1$ a linear, $a_2$ a quadratic form,
and so forth, and the corresponding holds of the coefficients
$b_0, b_1 \ldots b_m$ of the function $\varphi(x)$.

If $x, y, z$ denote triangular coordinates in the plane, then
$f(x) = 0$, $\varphi(x) = 0$ are the equations of a curve of order $n$ and
one of order $m$, and the common roots of these two
equations (the ratios $x:y:z$) denote the points of intersection
of the two curves. The final equation, which one obtains by
eliminating $x$, is a homogeneous equation whose degree
comes out, by (15), \S.\,48, equal to $nm$.

A lowering of the degree cannot take place here; the
final equation is always a homogeneous equation of degree $nm$
in the two unknowns $y, z$. The only exception to which
attention must be paid here is that the resultant vanishes
identically (for all $y, z$), so that infinitely many common
roots are present. In geometry this case has the meaning
that the two curves have a portion of a curve in common, and
not merely isolated points.

If the equations $f = 0$, $\varphi = 0$ have only a finite
number of common roots, then in the expressions

\[
y_1 = y + \beta x, \qquad z_1 = z + \gamma x
\]

we can arrange the constants $\beta, \gamma$ so that, whenever to one value
of the ratio $y : z$ there correspond several distinct values of $x$
satisfying the equations (1), the associated values
of $y_1 : z_1$ turn out different from one another.
